\documentclass[11pt]{article}
\usepackage[margin=1in]{geometry}
\usepackage{amsmath,amssymb}
\usepackage{hyperref}

\title{Discovering Hidden Relations Among Triangle Invariants}
\author{triangle-relations}
\date{}

\newcommand{\M}{\mathcal{M}}

\begin{document}
\maketitle

\section{Setup: the moduli space of triangles}

A triangle is specified by three points in $\mathbb{R}^2$, i.e.\ six real
parameters. Every quantity we consider (area, perimeter, circumradius $R$,
inradius $r$, distances between derived points, \dots) is built purely from
the vertex coordinates and is invariant under rigid motions: translating or
rotating the triangle does not change its area, its inradius, or the
distance between its incenter and circumcenter. Translations (2 parameters)
and rotations (1 parameter) form a 3-dimensional group acting on the
6-dimensional space of vertex triples; quotienting it out leaves a
$3$-dimensional moduli space $\M$ of triangle shapes up to congruence
(reflection is a further discrete symmetry and does not change this
dimension count). Concretely, $\M$ can be coordinatized by, e.g., the three
side lengths $(a,b,c)$ subject to the triangle inequality.

Every derived scalar quantity is a smooth function $f : \M \to \mathbb{R}$.
Choosing $k$ such quantities gives a smooth map
\[
  F = (f_1, \dots, f_k) : \M \to \mathbb{R}^k .
\]

\section{Why relations among four or more quantities are not interesting}

If $k > \dim \M = 3$, a generic smooth map from a 3-manifold into
$\mathbb{R}^k$ has, as its image, an (at most) 3-dimensional submanifold of
$\mathbb{R}^k$ — this is just the statement that embedding a 3-dimensional
space into a higher-dimensional ambient space produces $k-3$ constraint
equations for free, regardless of which $k$ functions were chosen, as long
as they are not specially degenerate. Concretely, this image is cut out
locally by $k-3$ independent equations $G_i(f_1,\dots,f_k) = 0$. So finding
\emph{some} relation among four or more triangle invariants is guaranteed by
dimension counting alone. It says nothing specific about triangle geometry —
it would happen for essentially any four smooth functions on any
3-dimensional space of objects.

\section{Why a relation among exactly three quantities is special}

If $k = \dim \M = 3$ exactly, the situation is qualitatively different. A
generic smooth map $F : \M \to \mathbb{R}^3$ is, at a generic point, a local
diffeomorphism: its Jacobian $DF$ has full rank 3, and as the triangle shape
ranges over an open subset of $\M$, the triple $(f_1,f_2,f_3)$ sweeps out a
full 3-dimensional \emph{open region} of $\mathbb{R}^3$ — a solid blob, not a
surface. This is again a genericity statement: for an arbitrary hand-picked
triple of triangle invariants, there is no a priori reason for $DF$ to be
rank-deficient anywhere.

An \emph{exact} relation $G(f_1,f_2,f_3) = 0$ holding on all of $\M$ is
therefore equivalent to $DF$ having rank $\le 2$ everywhere: the image
collapses onto a 2-dimensional surface embedded in $\mathbb{R}^3$, instead of
filling out a full 3-dimensional region. That is not generic — it is a
genuine algebraic coincidence, exactly the kind of fact that constitutes a
classical theorem. Euler's relation
\[
  d^2 = R^2 - 2Rr
\]
(where $d$ is the distance between the incenter and circumcenter) is exactly
a statement that the Jacobian of $(R, r, d)$ as functions of triangle shape
has rank $\le 2$ everywhere, collapsing three ostensibly independent
quantities onto a single surface. This is the precise sense in which
searching for relations among exactly three derived scalars — and not four —
is the search for something structurally new.

\section{Detecting rank-deficiency with a bottleneck autoencoder}

Given a candidate triple $(f_1,f_2,f_3)$, we want a numerical test for
whether the induced map $F : \M \to \mathbb{R}^3$ has full rank 3 generically
(data fills a 3D region) or is everywhere rank $\le 2$ (data is confined to a
2D surface). We sample many random triangles — many points of $\M$ — and
evaluate $(f_1,f_2,f_3)$ on each, producing a point cloud
$X \subset \mathbb{R}^3$.

An autoencoder with a 2-dimensional latent bottleneck, an encoder
$e : \mathbb{R}^3 \to \mathbb{R}^2$ and decoder $d : \mathbb{R}^2 \to
\mathbb{R}^3$ trained to minimize $\mathbb{E}\, \| d(e(x)) - x \|^2$ over the
sample, is exactly searching for the best rank-2 (nonlinear) approximation to
the data:
\begin{itemize}
  \item If the data truly lies on a 2-dimensional surface (an exact relation
    holds), a 2-dimensional latent code is not actually a bottleneck for it —
    a sufficiently expressive encoder/decoder pair can in principle drive the
    reconstruction error to (near) zero, since the encoder just needs to
    learn a local chart of the surface and the decoder its inverse.
  \item If the data genuinely fills an open 3-dimensional region, no
    continuous map that factors through a 2-dimensional intermediate space
    can reconstruct it without loss: an entire dimension worth of variation
    is necessarily discarded. Reconstruction error is then bounded away from
    zero, on the order of the data's own extent in the discarded direction.
\end{itemize}

This gives a clean signature: near-zero reconstruction error with a
2-dimensional bottleneck is evidence for an (exact or near-exact) relation
among the three chosen quantities; error comparable to the data's own scale
is consistent with the generic, full-rank, no-relation case.

\section{Calibrating ``near zero'' with a permutation null}

An absolute error threshold is not meaningful across different triples of
scalars: they differ in units, scale, and marginal shape, and — since all
derived scalars are ultimately functions of the same three underlying
triangle parameters — some triples are mildly correlated even without an
exact relation. What is needed is a triple-specific baseline for ``no
relation.''

We construct a null sample $X'$ by independently permuting each of the three
columns of $X$. This exactly preserves each $f_i$'s own one-dimensional
marginal distribution while completely destroying the joint dependency
between the three: each row of $X'$ combines values from unrelated triangle
instances. Any real functional relation is annihilated by this shuffle.
Training the same autoencoder architecture on $X'$ (repeated a few times)
estimates the reconstruction error one should expect for three quantities
with these particular marginals but \emph{no} shared structure. Comparing the
real error against this null's mean and standard deviation via a $z$-score
(or the ratio of the two) yields a self-calibrated significance measure that
is robust to the units and distribution of whichever three scalars are being
tested. A small ratio — real error far below the null — is the fingerprint
of an unexpected relation among exactly three quantities.

\section{From detection to an explicit formula}

The autoencoder step is a screening tool: it flags which triples are
suspicious, not what the relation says. Once a triple is flagged, we look for
an explicit closed form by fitting a polynomial identity. We build a
dictionary of monomials in $(f_1,f_2,f_3)$ up to some degree $D$ (for $D=2$:
$1, f_1, f_2, f_3, f_1^2, f_1 f_2, f_1 f_3, f_2^2, f_2 f_3, f_3^2$), evaluate
them on the sampled data to form a design matrix $\Phi$ (samples $\times$
monomials, each column rescaled to comparable variance), and look for a
linear combination of its columns that vanishes on every sample. Such a
combination is the right singular vector belonging to the smallest singular
value of $\Phi$.

Applied to $(R, r, d)$ at degree 2, this recovers — up to overall scale and
sign, exactly the ambiguity inherent in reading off a null vector —
\[
  d^2 - R^2 + 2Rr \;\approx\; 0
\]
on every sampled triangle: Euler's relation, found purely by sampling random
triangles and asking a null-space question, with no prior knowledge of the
formula. A very small trailing singular value (relative to the largest) is
itself independent confirmation that an exact low-degree polynomial relation
exists, complementary to the autoencoder's nonlinear, degree-agnostic
detector.

\section{The full pipeline}

\begin{enumerate}
  \item Enumerate candidate triples of derived scalar quantities.
  \item For each triple, sample random triangles once, train a 2-bottleneck
    autoencoder, and compare its held-out reconstruction error to a
    column-shuffled null (several repeats).
  \item Rank triples by how much better the real data compresses than its
    null counterpart (small ratio $=$ strong candidate).
  \item For top candidates, fit a low-degree polynomial null-space relation
    to recover an explicit closed form, and validate that it holds (residual
    $\approx 0$) on fresh samples.
\end{enumerate}

\texttt{triangle\_relations/discovery/verify\_euler\_relation.py} runs this
exact pipeline on $(R, r, d)$ as a worked example, confirming it rediscovers
Euler's relation end to end.

\end{document}
